run\documentclass[11pt]{article}

\usepackage[margin=1in]{geometry}
\usepackage{amsmath,amssymb}
\usepackage{booktabs}
\usepackage{tabularx}
\usepackage{hyperref}

\setlength{\parindent}{0pt}
\setlength{\parskip}{0.6\baselineskip}

\title{Recurrent Offline Reinforcement Learning for OGBench CRL\\
\large Complete Project Specification}
\author{}
\date{Last updated: February 7, 2026}

\begin{document}
\maketitle


\begin{abstract}
Offline goal-conditioned reinforcement learning (GCRL) promises to synthesize long-horizon behaviors from static datasets, but performance remains limited by generalization. We study whether \emph{iterated, weight-tied} critic backbones improve compositional generalization in offline GCRL while preserving parameter efficiency. Using OGBench and the standard Contrastive RL (CRL) pipeline, we replace only the contrastive critic encoders ($\phi(s,a)$ and $\psi(g)$) with an iterative residual backbone applied for $K$ steps, using step conditioning (FiLM) and stabilizers (LayerNorm and gated/LayerScale residuals). This modification is architecture-only: the CRL objective, dataset, actor, training budget, ensemble semantics, and evaluation protocol are kept unchanged. On \texttt{antmaze-large-stitch-v0}, the best tied iterative critic (sinusoidal, $1\times4$) achieves $0.475 \pm 0.039$ success across three seeds, substantially exceeding feedforward ResNet baselines ($0.257 \pm 0.025$ for depth 4; $0.203 \pm 0.080$ for depth 8). Iterations beyond $K=4$ are not strictly monotonic, suggesting an optimization/overfitting tradeoff rather than ``more iterations always helps.'' Overall, our results indicate that weight-tied iterative critics are a strong and practical lever for improving offline goal-conditioned CRL on stitching tasks with favorable parameter efficiency.
\end{abstract}

\section{Introduction}
Offline goal-conditioned RL aims to learn policies that reach user-specified goals using only a fixed dataset of previously collected experience. In long-horizon navigation-like domains, offline learning is often bottlenecked not by low-level control but by \emph{compositional generalization}: to reach a novel goal, an agent must compose partial skills found in the dataset into a new plan. OGBench ``stitch'' tasks isolate this challenge by evaluating whether an agent can combine observed sub-trajectories into successful long-horizon behavior.

A major empirical difficulty in offline GCRL is critic generalization. In CRL, the critic is contrastive rather than a scalar $Q$-function: it learns representations via encoders $\phi(s,a)$ and $\psi(g)$, and training encourages alignment between state-action features and goal features. While the contrastive formulation can be effective, the representational backbone used for $\phi/\psi$ can still overfit and may fail to propagate value-relevant information across the manifold of states and goals needed for stitching.

Motivated by recent interest in recurrent/iterative computation for reasoning-like behaviors (i.e., performing multiple refinement steps in a latent space), we investigate a simple architectural hypothesis: \emph{replacing the feedforward critic backbone with an iterated weight-tied residual block} improves compositional generalization under a fixed training budget. Concretely, we keep CRL intact and only swap the critic backbone: a single residual MLP block is applied for $K$ iterations, with per-step conditioning (FiLM) and stabilization (LayerNorm plus gated/LayerScale residual updates). This design allows us to control effective depth via the iteration count while keeping parameters small, and it provides a clean testbed for whether iterative computation helps on stitching.

We evaluate on \texttt{antmaze-large-stitch-v0} with three seeds per configuration. Feedforward ResNet baselines (untied depth) achieve $0.257 \pm 0.025$ (depth 4) and $0.203 \pm 0.080$ (depth 8), while tied iterative critics significantly improve success: for example, a sinusoidal step-conditioned $1\times4$ model reaches $0.475 \pm 0.039$. Iteration scaling is not strictly monotonic ($1\times6$: $0.415 \pm 0.041$, $1\times8$: $0.455 \pm 0.052$), suggesting that unlocking test-time scaling will require additional training strategies (e.g., randomized-$K$ training) or improved step-conditioning. Nonetheless, the consistent gap between tied iterative critics and untied ResNets indicates that weight tying plus iterative refinement is a meaningful architectural lever for offline CRL on stitching tasks. \par

\section{Contributions}
\begin{itemize}
  \item \textbf{Architecture-only improvement to offline CRL.} We present a drop-in replacement for CRL's contrastive critic backbone (the $\phi(s,a)$ and $\psi(g)$ encoders), using an iterated weight-tied residual block with step conditioning and stabilization, while leaving the CRL objective, actor, training loop, ensemble semantics, and evaluation unchanged.

  \item \textbf{Strong gains on AntMaze stitching under fixed budget.} On \texttt{antmaze-large-stitch-v0}, tied iterative critics substantially outperform feedforward ResNet baselines under the same training budget and comparable widths. For example, sinusoidal $1\times4$ achieves $0.475 \pm 0.039$ vs.\ ResNet depth-4 at $0.257 \pm 0.025$ (3 seeds).

  \item \textbf{Depth/iteration decomposition for controlled study.} We frame critic compute as a decomposition into \emph{stack depth} and \emph{outer iterations} (e.g., $1\times K$, $2\times2$, $4\times1$), enabling systematic comparisons among feedforward depth, partial tying (UT-style), and fully tied recurrence.

  \item \textbf{Evidence of non-monotonic iteration scaling.} Increasing $K$ beyond 4 is not strictly monotonic in current settings, highlighting that ``test-time compute'' is not automatically unlocked by recurrence and motivating future methods such as randomized-$K$ training or improved step-conditioning.
\end{itemize}


\begin{table}[b]
\centering
\small
\begin{tabular}{l l c l c}
\hline
\textbf{Model} & \textbf{Backbone} & \textbf{Iter} & \textbf{Seeds} & \textbf{Success} \\
\hline
ResNet $d=4$ & Feedforward & Fixed & 0.23, 0.29, 0.25 & $0.257 \pm 0.025$ \\
ResNet $d=8$ & Feedforward & Fixed & 0.30, 0.10, 0.21 & $0.203 \pm 0.080$ \\
Discrete $1\times4$ & Recurrent (discrete) & 4 & 0.46, 0.41, 0.36 & $0.410 \pm 0.041$ \\
Tied $2\times1$ & Recurrent (sinusoidal) & 2 & 0.344, 0.26, 0.124 & $0.243 \pm 0.091$ \\
Tied $4\times1$ & Recurrent (sinusoidal) & 4 & 0.28, 0.272, 0.212 & $0.255 \pm 0.030$ \\
Tied $1\times2$ & Recurrent (sinusoidal) & 2 & 0.332, 0.212, 0.128 & $0.224 \pm 0.084$ \\
Tied $1\times6$ & Recurrent (sinusoidal) & 6 & 0.364, 0.416, 0.464 & $0.415 \pm 0.041$ \\
Tied $1\times8$ & Recurrent (sinusoidal) & 8 & 0.384, 0.508, 0.472 & $0.455 \pm 0.052$ \\
Tied $1\times4$ & Recurrent (sinusoidal) & 4 & 0.436, 0.528, 0.460 & $0.475 \pm 0.039$ \\
\hline
\end{tabular}
\caption{Current results on \texttt{antmaze-large-stitch-v0} (3 seeds each).}
\label{tab:antmaze_large_stitch}
\end{table}


\section{Reinforcement Learning Algorithm}
\subsection{Base algorithm: CRL (Contrastive RL)}
\textbf{Q-value definition:}
\[
  Q(s,a,g) = \mathbb{E}\left[\sum_{t=0}^{\infty} \gamma^t r_t\,\middle|\, s_0=s, a_0=a, \text{goal}=g\right].
\]
In CRL, the critic is represented as a bilinear score
\[
  Q(s,a,g) = \frac{\phi(s,a)^\top \psi(g)}{\sqrt{d}}.
\]

\subsection{Key components}
\subsubsection{Contrastive loss (sigmoid BCE)}
The critic is trained to discriminate positive triplets $(s,a,g)$ from negatives using a sigmoid binary cross-entropy objective on logits $\phi^\top\psi/\sqrt{d}$ with labels $\mathbb{I}$. Positive pairs correspond to observed state--action and achieved-goal pairings; negatives are produced by mismatching goals within the same batch (in-batch negatives). Empirically, batch size $B=1024$ is sufficient.

\subsubsection{Bilinear critic architecture}
\begin{itemize}
  \item \textbf{State-action encoder:} $\phi : (s,a) \mapsto \mathbb{R}^d$.
  \item \textbf{Goal encoder:} $\psi : g \mapsto \mathbb{R}^d$.
  \item \textbf{Scaling:} $1/\sqrt{d}$.
  \item \textbf{Ensemble:} two critics; the minimum is used for robustness.
\end{itemize}

\subsubsection{Actor training}
The actor is trained with DDPG+BC using
\[
  \mathcal{L}_{\text{actor}} = \alpha\,\mathrm{MSE}(\pi(s,g), a_{\text{data}}) - Q\bigl(s,\pi(s,g),g\bigr),
\]
with behavior cloning coefficient $\alpha=0.1$. The policy $\pi:(s,g)\mapsto a$ is deterministic and outputs actions clipped to $[-1,1]$.

\section{Neural Network Architectures}
\subsection{Baseline MLP critic (untied)}
\textbf{Architecture (latent dimension $512$):}
\begin{itemize}
  \item $\phi(s,a)$: concatenate $[s;a]$ then apply four dense layers of width 512.
  \item $\psi(g)$: apply four dense layers of width 512.
\end{itemize}
Layer normalization is enabled (applied after each layer). Each backbone has approximately 1.3M parameters.

\subsection{Deep ResNet critic (untied)}
\textbf{Structure:} $[s;a] \to \mathrm{Dense}(512) \to [\mathrm{ResBlock}]^{\times D} \to \mathrm{Dense}(512)$ with $D\in\{4,8\}$.

\textbf{Residual block (pre-norm):}
\[
  h_1 = \mathrm{LN}(h),\quad
  u = \mathrm{Dense}(512)\bigl(h_1\bigr),\quad
  u = \mathrm{SiLU}(u),\quad
  u = \mathrm{Dense}(512)(u),\quad
  h \leftarrow h + \mathrm{layerscale}\cdot u.
\]
\textbf{Known issues:} depth $D=8$ exhibits Q-value collapse (e.g., Q-mean drift from $-4.0$ to $-6.0$), increased contrastive loss, and worse final performance than $D=4$. Standard feedforward ResNets also (i) lack explicit goal-conditioning (no FiLM modulation) and (ii) cannot scale depth at test time.

\subsection{Recurrent tied critic (proposed)}
\textbf{Core idea:} weight-tied iterative refinement with FiLM conditioning.
\begin{itemize}
  \item Backbone: $[s;a] \to \mathrm{Dense}(h) \to [\mathrm{TiedBlock}]^{\times K} \to \mathrm{Dense}(d)$.
  \item Training iterations: $K_{\text{train}}=4$.
  \item Test iterations: $K_{\text{test}}\in\{4,8,12,16\}$ (test-time compute scaling).
  \item Hidden dimension: $h=512$; latent dimension: $d=512$.
  \item Per-step scaling: learned $\alpha_k$ (LayerScale), initialized at $10^{-2}$.
\end{itemize}

\textbf{Tied iteration block (high-level):}
\begin{itemize}
  \item Shared LayerNorm.
  \item Dynamic FiLM conditioning: per-iteration $(\gamma,\beta)$ depends on both the step index and the current hidden state.
  \item Efficiency: split-fc1 precomputation, where $\mathrm{fc1}([\text{step};h_1])$ is decomposed as \mbox{$\mathrm{fc1}_{\text{step}}(\text{step}) + \mathrm{fc1}_{h}(h_1)$}.
\end{itemize}


\subsection{Actor network}
\textbf{Architecture:} $[s;g] \to \mathrm{Dense}(512) \to \mathrm{Dense}(512) \to \mathrm{Dense}(512) \to \mathrm{Dense}(\mathrm{action\_dim})$.

\section{Positional Encoding for Test-Time Scaling}
\subsection{Discrete step embeddings (baseline)}
A learnable table $\mathrm{Embed}(\mathrm{max\_iters}, h)$, with $\mathrm{max\_iters}=16$ and $h=512$, totaling $16\times 512 = 8192$ parameters. Extrapolation beyond $K_{\text{train}}$ is poor because embeddings for $k > K_{\text{train}}$ are untrained.

\subsection{Sinusoidal positional encoding (proposed)}
A closed-form (non-learned) sinusoidal encoding enables smooth extrapolation to $K_{\text{test}} > K_{\text{train}}$ and supports test-time compute scaling.
\begin{itemize}
  \item Learnable parameters: 0.
  \item Frequency range: $[1.0, 10^{-4}]$.
  \item Empirical effect: scaling from $K=4$ to $K=8$ yields an absolute success improvement of $0.038$ (from $0.475$ to $0.513$).
\end{itemize}

\section{Training Hyperparameters}
\subsection{Optimization}
\begin{itemize}
  \item Optimizer: Adam.
  \item Learning rate: $3\times 10^{-4}$ with cosine decay to $10^{-5}$ over 1,000,000 steps (recurrent-only).
  \item Batch size: 1024.
  \item Discount factor: $\gamma=0.99$ (and $0.995$ for giant).
  \item Gradient clipping: enabled (via global norm clipping; logged).
\end{itemize}


\subsection{CRL-specific hyperparameters}
\begin{itemize}
  \item Behavior cloning coefficient: $\alpha=0.1$.
  \item Actor goal sampling: $p_{\text{random}}=0.5$, $p_{\text{traj}}=0.5$.
  \item Critic goal sampling: $p_{\text{traj}}=1.0$; future goals sampled geometrically.
  \item Ensemble size: 2.
\end{itemize}



\section{Future Work}
\begin{enumerate}
  \item Extend evaluation to additional OGBench tasks (e.g., pointmaze).
  \item Train anytime predictors with randomized $K \in [1, K_{\max}]$ and/or early-exit mechanisms.
  \item Investigate hybrid architectures combining partial tying with sinusoidal encoding.
\end{enumerate}

\end{document}